\section{Introduction}\label{sec:intro}
Many-agent systems need a reliable answer to ``which declared checks hold on
this exact candidate?'' Their task graph and conversational organization need
not be the layout that aggregates this evidence. HACT separates the two:
immutable check identities determine the acceptance contract, while a
bounded-packet tree represents intermediate results. A smaller tree transcript
is not automatically a faster build, a better agent, or a smaller root-status
prompt.

The distinction matters economically. Historical source timings were 39.69
versus 39.77 seconds for fixed and learned layouts. A separate small harness
paid roughly five seconds per episode for strict verification compared with
direct pytest. Meanwhile, a historical compressed loopback export saved
approximately 2.21 seconds across twelve source streams. These different
cohorts cannot be combined into a measured end-to-end speedup, and neither
observation justifies adding a strict wrapper solely to save transmission time.

\paragraph{Three deployment decisions, not one.}
A trusted root-status consumer needs neither a learned tree nor all intermediate
packets. A deployment that already requires protected verification and
hierarchical evidence can compare fixed and learned layouts at marginal cost
under that unchanged service. Adopting protected verification where ordinary
pytest was sufficient is a different decision: its additional preparation and
checking costs must be justified independently. We make this service distinction
explicit in both the paper and a fail-closed planning API.

\paragraph{The simpler compiler is now primary.}
The prior block compiler was 1.54--5.34\% more expensive in held-out bytes than
a compact balanced eight-ary tree on the same learned order. We therefore remove
block DP from the practical path. The new default catalogue retains the
incumbent and proposes a balanced tree on a training-learned average-linkage
order; separate validation compares full-wave costs. Exact DP remains a useful
conditional optimum and reference. It can still produce smaller evidence than the
practical default, as our source replays show. Low fitting cost and lowest byte
cost are different objectives.

\paragraph{Core contributions and architecture.}
Adaptive HACT v5.1 serializes verification evidence under fixed test execution
schedules. Its architecture separates two kinds of state that many systems
conflate. \emph{Immutable check identities}---registry IDs, acceptance
predicates, and the candidate binding---define what must hold for a candidate
to be accepted; they are fixed when the registry and candidate identity are
published and are never rewritten by execution, layout choice, or learning.
\emph{Mutable execution state}---tokens, timestamps, generation counters, and
the order in which checks actually ran---records one run of those checks and
may differ run to run. A result token carries the immutable semantic ID plus
the run's binding; a layout-specific kernel maps stable IDs to display
positions. Because identities are stable, permuting the layout or replaying
under a different schedule moves records without changing their meaning
(Section~\ref{sec:soundness}). The execution schedule $\sigma_e$ is an input
the serializer must respect, not a variable it optimizes: HACT changes how
evidence from an unchanged schedule is laid out and aggregated, never when
checks run or what their outcomes mean. The substantive research contribution
remains hyperedge-scored, identity-preserving verification aggregation and its
conditional guarantees. This revision resolves an implementation choice, adds
service-matched cost accounting, and evaluates the no-DP practical catalogue
on fresh frozen synthetic splits. It separately replays archived source
evidence and provides tools for the unmeasured token and WAN boundaries.
Clustering, balanced trees, alphabetic DP, compression, break-even arithmetic
and TLS are established techniques, not new inventions claimed here. The
artifact is a tested research implementation; its tests do not establish
production readiness or top-conference acceptance.

\paragraph{Positioning: evidence serialization, not agent acceleration.}
Every design decision in v5.1 follows from a deliberately narrow scope: HACT is
an evidence-serialization component. It reduces the bytes needed to publish
which declared checks hold, under a delivery service that already requires
protected verification and hierarchical evidence. It does not schedule checks,
does not parallelize work, does not alter the agent's planning or repository
exploration, and does not claim that a smaller transcript makes any agent
faster or more productive. Where the review's larger program---verified-frontier
context retirement, dependency-aware parallelism, task granularity---would use
this certificate as a backend, that program is an evaluation protocol in
progress (Section~\ref{sec:productivity}); its experiments have not been run.
The negative controls in Section~\ref{sec:productivity} exist precisely to
keep this boundary honest: flat checker timings and per-episode verification
overhead are reported as limits of the serialization claim, not evidence
against it or support for an acceleration claim.
